% RESEARCH_DRAFT_SHA256: 9ac4962effda4e137a76f1c6d4e30c9c9b2e0e23ac60fc25f77a563bb3c77e8f
% Layout and source-reference wrappers expanded for the web renderer.
\documentclass[11pt]{article}
\usepackage[margin=25mm]{geometry}
\usepackage{amsthm}
\usepackage{fontspec,unicode-math}
\setmainfont{Latin Modern Roman}
\setmathfont{Latin Modern Math}
\usepackage{amsmath,amssymb,mathtools,xurl,hyperref}
\hypersetup{colorlinks=true,urlcolor=blue,linkcolor=blue}
\setcounter{secnumdepth}{0}
\setlength{\parindent}{0pt}
\setlength{\parskip}{5pt}
\setlength{\emergencystretch}{3em}
\newcommand{\sref}[2]{\hyperlink{source:#1:#2}{[S#2]}}
\newcommand{\cataloguescope}[1]{\paragraph{Source record.}#1}
\newcommand{\floor}[1]{\left\lfloor#1\right\rfloor}
\newtheorem{proposition}{Proposition}
\newtheorem{lemma}[proposition]{Lemma}
\begin{document}
\paragraph{Research manuscript.} OpenAI \texttt{gpt-6-astra}, reasoning \texttt{ultra}; 11 October 2026.
\paragraph{Preserved source body.} The complete catalogue mathematical text and sources follow with layout and source links expanded, including source qualifications and its historical preparation notes. New research begins at the section headed Research attempt.

% UnsolvedMath ID 30006663; source catalogue code OWR-14299916-001
\section{Complexity and Generalization in Overparameterized Learning}\label{30006663}
{\small UnsolvedMath ID 30006663 · Computer Science}
\subsection{Definitions and mathematical statement}

\paragraph{Shared notation.}Standard number systems are $\mathbb N,\mathbb Z,\mathbb Q,\mathbb R,\mathbb C$. Specialized notation and index ranges follow the cited source and the definitions below.


Overparameterization compares the number of trainable parameters with the stated data or effective-complexity scale. Implicit bias is the preference of an optimization procedure among solutions with comparable training fit. Population loss measures generalization under the target distribution.
\paragraph{Record scope and qualifications.}
The catalogue title is retained as metadata. Named certificates, contracts, bounds or frameworks in that title are not asserted as results of the original source extract.
These are broad theory directions; the architecture, loss, data law, and requested bounds are not fixed.

Source-record limitation: The listed broad learning-theory directions do not specify mathematical propositions or quantitative targets.
\paragraph{Statement recorded in the supplied source.}
Which intrinsic notions of complexity govern learning and generalization in highly and moderately overparameterized models? Why do gradient-based optimization methods consistently select solutions that generalize well? How do useful representations emerge during training, and what are the computational limitations of modern learning systems? \href{https://doi.org/10.4171/owr/2026/11}{[S2]}
\paragraph{Source status.}
The supplied catalogue labels this source record open; that metadata was checked against the supplied snapshot on 2026-10-11. The label is preserved for provenance and is not an independent verification of present open status or mathematical correctness. The catalogue literature assessment is dated 2026-08-22.
\subsection{Short English statement}
Explain complexity, optimization bias, representation learning, and computational limits in modern learning.
\subsection{Sources}
\begin{itemize}
\item[S1] Ulam AI and the UnsolvedMath Contributors, UnsolvedMath: A Curated Collection of Open Mathematics Problems, record 30006663 (OWR-14299916-001). Catalogue curation and metadata: CC BY 4.0. The source extract is retained; definitions, English summaries and qualifications are editorial additions.
\url{https://huggingface.co/datasets/ulamai/UnsolvedMath}.
\item[S2] Modern and Emerging Phenomena in Machine Learning (2026), Oberwolfach Reports, DOI 10.4171/owr/2026/11, pp. 683--736
\url{https://doi.org/10.4171/owr/2026/11}.
\item Source-specific Oberwolfach Reports licences, including CC BY-SA 4.0 where applicable, are retained; see THIRD\_PARTY\_NOTICES.md.
\end{itemize}
\paragraph{Preparation and review.}
The original source and UnsolvedMath contributors retain attribution for the source extract. Codex (GPT-6; exact serving revision not exposed) assisted with definitions, English summaries, scope notes and formatting. No new human mathematical authorship or independently verified solution is asserted. Permanent source ID 30006663 and code OWR-14299916-001 checked against the complete supplied registry and current repository. Distinct upstream records, including repeated questions and broad or damaged extracts, are retained for complete coverage of the source CS category. A repeated record is not a new independent conjecture. Complete mathematical-equivalence and current-literature review has not been performed. Preparation and final review follow contribution prompts C and D, with source gaps disclosed. Structural and rendering checks do not establish a new theorem or open status.
\label{end:30006663}

\section{Research attempt}
\subsection{Literature, source context and scope}
The full catalogue body, including qualifications and historical source status, is preserved above. The original workshop report was inspected in the locally cached publisher PDF text. This attempt distinguishes the source question from any explicitly restricted model introduced below. Repeated records retain their own source IDs and do not become independent conjectures. Elementary results and attributed literature results are not claimed as new original-target progress.

\subsection{Testing an unqualified implicit-bias explanation}
The source poses broad directions without a fixed architecture or data law. We test the inference that the minimum-norm solution selected by gradient flow necessarily generalizes well. The following deliberately specified linear model isolates what the training data fail to identify.
There are two parameters $\theta=(\theta_1,\theta_2)$ and one training example $x=(1,0)$ with label one. The squared training loss is $L(\theta)=\tfrac12(\theta_1-1)^2$. Start Euclidean gradient flow at zero. For any $M\in\mathbb R$, let the true response on a test vector $X$ be $Y=X_1+MX_2$, with $\mathbb E[X_2^2]=1$.
\begin{proposition}
Gradient flow converges to the unique minimum-Euclidean-norm interpolant $(1,0)$, whose test mean squared error is $M^2$.
\end{proposition}
\begin{proof}
The equations are $\dot\theta_1=1-\theta_1$ and $\dot\theta_2=0$, giving $\theta(t)=(1-e^{-t},0)$. All interpolants have the form $(1,a)$, whose squared norm $1+a^2$ is minimized uniquely at $a=0$. At this limit, prediction minus truth is $-MX_2$, and its squared expectation is $M^2$. Every value of $M$ produces exactly the same given training example, since its second feature is zero.
\end{proof}
Thus neither interpolation nor the optimizer's minimum-norm bias supplies a distribution-free generalization guarantee. One needs information about the target's unseen directions, the sampling law or an appropriate relation between norm and test covariance. This does not refute specific learning theorems with those assumptions, and this fixed training design is not being represented as a typical random sample.
\subsection{Verification and remaining gaps}
At $M=10$ and test vectors $X=(0,1),(0,-1)$ with equal probabilities, both squared errors are 100 while limiting training error is zero. No intrinsic complexity measure for modern networks, representation-emergence theory, or computational-limit characterization is produced.

\section{References and source attribution}
The primary Oberwolfach workshop report, \url{https://doi.org/10.4171/owr/2026/11}, was inspected on 11 October 2026. Relevant source passages: 2026/11 pp. 684--685, introductory complexity and generalization questions. Catalogue and contributor attribution remain in the original Sources section. Mathematical arguments labeled as this attempt are proved above; cited external results remain attributed to their authors and are not independently reproved or claimed as our discoveries.


\paragraph{Integration review.} Codex (GPT-6; exact serving revision not exposed) reviewed the written arguments and expanded layout and source-reference wrappers. Repository metadata and the production web renderer were checked using Python 3.13 and JavaScript with MathJax 3.2.2. This model-assisted review is not independent human proof verification; no PDF compilation or formal proof checking is asserted.

\section{Completion Estimate}
\noindent\textbf{COMPLETION: 0\%}
No substantiated new progress toward the entire original target is claimed. The worked results and failed route are delimited above; a known result, illustrative model, or auxiliary obstruction is not counted as original-target completion.

\section{Final status and attribution}
\textbf{Attempt status: unresolved.} This describes this research attempt, not a new ruling on the source catalogue's status or on external literature claims.
The mathematics and this self-audit were generated by OpenAI \texttt{gpt-6-astra}, reasoning \texttt{ultra}, on 11 October 2026. Codex (GPT-6; exact serving revision not exposed) orchestrated the batch and contributed formatting/repository review. No human mathematical authorship, independent human proof checking, or formal verification is claimed. The historical preparation paragraphs in the reproduced catalogue body concern that body only.

Contribution prompt A was applied to the complete definition and supplied source record, followed by the adversarial checks of prompt B. Local reproduction files are \texttt{batch011/checks.py}, \texttt{computations.json}, and the per-ID audit JSON. They were not public links at preparation time. Numerical checks are finite and do not prove unrestricted results. No LaTeX engine was found locally; compilation was not performed by this research agent.

\end{document}
